% !PRINTED 267
\weberpage{267}
%<<CONT
setzen, so dargestellt werden kann:
\weq{9}{\Phi_p(v,u)=\Bigl(\frac{uv}{\sqrt2}\Bigr)^{p+1}\Phi_p\Bigl(\frac{\varepsilon\sqrt2}{v},\frac{\sqrt2}{u}\Bigr).}
Hieraus schlie\ss t man auf die Relationen
\weq{10}{\varepsilon^h2^{\frac{h+k-p-1}{2}}a_{p+1-h,\,p+1-k}=a_{h,k}, \qquad a_{h,k}=a_{k,h}.}
Wenn wir also in (7) die Glieder mit gleichen Koeffizienten
$a_{h,k}$ zusammenfassen, so k\"onnen wir uns auf die Annahme
beschr\"anken, da\ss\ $h\geqq k$, $h+k\geqq p+1$ sei, und es ergibt
sich $\Phi_p$, von den beiden Gliedern $v^{p+1}, u^{p+1}$ abgesehen,
als ein Aggregat von Gliedern von den folgenden beiden Formen
(wenn noch ber\"ucksichtigt wird, da\ss\ wegen (8) $\varepsilon^h=\varepsilon^k$ ist):
\weq{11}{\begin{aligned}
&v^hu^k+v^ku^h+\varepsilon^h2^{\frac{h+k-p-1}{2}}\bigl(u^{p+1-h}v^{p+1-k}+u^{p+1-k}v^{p+1-h}\bigr)=\\
&\underset{h>k}{(uv)^{\frac{p+1}{2}}}\left[\Bigl(\frac{v}{u}\Bigr)^{\frac{h-k}{2}}+\Bigl(\frac{u}{v}\Bigr)^{\frac{h-k}{2}}\right]\left[(uv)^{\frac{h+k-p-1}{2}}+\frac{\varepsilon^h2^{\frac{h+k-p-1}{2}}}{(uv)^{\frac{h+k-p-1}{2}}}\right].
\end{aligned}}
\weq{12}{\begin{aligned}
&v^hu^h+\varepsilon^h2^{h-\frac{p+1}{2}}(uv)^{p+1-h}\\
&=(uv)^{\frac{p+1}{2}}\left[(uv)^{h-\frac{p+1}{2}}+\frac{\varepsilon^h2^{h-\frac{p+1}{2}}}{(uv)^{h-\frac{p+1}{2}}}\right].
\end{aligned}}
Die Koeffizienten dieser Glieder in $\Phi_p$ sind ganze Zahlen.
Nun ergibt sich aus (8):
\[
\begin{aligned}
(h-k)&\equiv (h-1)(p+1)\\
h+k-p-1&\equiv -h(p-1)
\end{aligned}
\ (\mathrm{mod}\ 24),
\]
und daher ist $h-k$ durch $2r$, $h+k-p-1$ (worin $h$ auch
$=k$ sein kann) durch $2s$ teilbar [\S\ 74, (2)].

Setzen wir daher
\[
\frac{h-k}{2}=r\alpha, \qquad \frac{h+k-p-1}{2}=s\beta,
\]
so sind $\alpha$ und $\beta$ positive ganze Zahlen, die zwischen
den Grenzen
\weq{13}{0\leqq\alpha<\frac{p+1}{2r}, \qquad 0\leqq\beta\leqq\frac{p+1}{2s},}
liegen, und \"uberdies ergibt die aus (8) folgende Kongruenz
$-h(p-1)\equiv 2s\beta\ (\mathrm{mod}\ 24)$, da\ss\ wenigstens in den
F\"allen, wo $\varepsilon=-1$ ist, $h\equiv\beta\ (\mathrm{mod}\ 2)$
[\S\ 74, (2)], also stets $\varepsilon^h=\varepsilon^\beta$.
